\section{Introduction}
\label{sec:introduction}

Two questions about anticanonical bundles motivate this paper. First,
when a torus acts on a compact complex manifold, can the invariant Euler
characteristic at the trivial power force invariant cohomology at large
anticanonical powers? Second, on a smooth projective variety, can a fixed
pseudoeffective error be removed from an unbounded sequence of effective
anticanonical twists? We answer the first question without a positivity
assumption and the second under smooth metric semipositivity. They meet
in an application to differential forms: invariant cohomology supplies
forms that descend through monodromy, and conversion is applied only
after that descent.

\subsection{The invariant index and its value at zero}

Let \(T\) be a compact real torus acting smoothly by holomorphic
automorphisms on a compact complex manifold \(F\). Its differential
acts on the anticanonical line \(L_F=-K_F=\det T_F\). We call this
the \emph{natural linearization}. The corresponding action on a section
is \((t\cdot s)(x)=t_{L_F}s(t^{-1}x)\); tensor powers carry the induced
action. For each integer \(m\ge0\), define
\[
I_T(m)=\sum_{q\ge0}(-1)^q\dim H^q(F,L_F^m)^T.
\]
Thus \(I_T(m)\) is the multiplicity of the trivial representation in the
Dolbeault index of \(L_F^m\).

\begin{theorem}[Invariant anticanonical index]\label{thm:index}
Let a compact real torus \(T\) act holomorphically on a compact complex
manifold \(F\), and give \(-K_F\) its natural linearization.
There are an integer \(a>0\) and a polynomial \(P\in\Q[t]\) such that
\[
I_T(m)=P(m)\quad\text{for every positive multiple \(m\) of \(a\)},
\qquad P(0)=I_T(0).
\]
\end{theorem}

The theorem has no projectivity, K\"ahlerness, or positivity premise.
Its endpoint is essential: when \(I_T(0)\ne0\), the polynomial cannot
vanish at all large divisible exponents. Some fixed cohomological degree
therefore contains invariant classes at unbounded exponents. Eventual
polynomiality alone would give no information about whether that
polynomial is zero.

The natural linearization is equally essential. On a point, twisting the
trivial line by a nonzero torus character gives invariant dimension zero
at every positive power and dimension one at power zero. For the natural
anticanonical action, by contrast, the fiber character at a fixed
component is the sum of the normal tangent characters. This identity
controls the endpoint.

Our calculation uses the holomorphic Lefschetz framework of
Atiyah--Bott \cite{AtiyahBott1966} and its compact-group, fixed-component
form due to Atiyah--Segal and Atiyah--Singer
\cite{AtiyahSegal1968,AtiyahSinger1968}. Polarizing its denominators
expresses the invariant multiplicity as polynomially weighted lattice
counts on rational polytopes with specified coordinate faces removed.
Weighted Ehrhart theory supplies polynomiality on divisible exponents
\cite[preprint \S2.4 and \S3.1]{BaldoniBerlineDeLoeraKoeppeVergne2012}.
To determine the value at zero, we prove the count directly, retaining
all small divisible exponents, and then show that the anticanonical
character gives a deformation retraction onto the removed faces.
The cancellation is a difference of Euler characteristics of compact
polyhedra. Stanley's treatment of deleted boundary already identifies
this constant term and singles out boundary visible from an exterior
point
\cite[Proposition~8.2 and the discussion and proof of Proposition~8.3]{Stanley1974Reciprocity}.
Here the anticanonical character produces such an exterior point in the
affine counting space. We give a direct weighted argument and an explicit
retraction, including degenerate polytopes.

When \(F\) is K\"ahler and \(L_F\) has a smooth semipositive metric,
hard Lefschetz with semipositive coefficients
\cite[Theorem~0.1]{DPS2001} gives
\[
H^0(F,\Omega_F^{n-q}\otimes L_F^{m+1})
\longrightarrow H^q(F,L_F^m),\qquad n=\dim F,
\]
surjectively. An invariant K\"ahler form makes this map equivariant;
compact averaging then preserves surjectivity on invariant subspaces.
The coefficient is \(L_F^{m+1}\) because the cohomological target is
\(K_F\otimes L_F^{m+1}=L_F^m\). Thus Theorem~\ref{thm:index} supplies
invariant twisted differential forms when \(I_T(0)\ne0\).
The use of hard Lefschetz and a nonzero Euler characteristic to obtain
infinitely many twisted forms already appears for canonical powers in
\cite[Theorem~2.7.3]{DPS2001}; here the invariant anticanonical index
requires a separate argument to retain its value at exponent zero.

\subsection{Removing a fixed error}

We use additive notation for line bundles and Cartier divisors. A divisor
is \emph{pseudoeffective} if its numerical class lies in the closure of
the effective cone. Smooth semipositivity means the existence of a smooth
Hermitian metric with nonnegative Chern curvature; it implies nefness.
The next theorem is an ordinary section theorem, with no group action.

\begin{theorem}[Pseudoeffective-error conversion]\label{thm:conversion}
Let \(S\) be a smooth connected projective complex variety, and suppose
that \(L=-K_S\) has a smooth Hermitian metric with semipositive Chern
curvature. Let \(D\) be a pseudoeffective Cartier divisor. If
\[
H^0(S,\OO_S(m_jL-D))\ne0
\]
for a strictly increasing sequence of positive integers \(m_j\), then
\(H^0(S,\OO_S(kL))\ne0\) for some integer \(k>0\).
\end{theorem}

The fixed error need not be effective. No rational-connectedness or
Euler-characteristic assumption on \(S\) is used. The argument begins
with a rational map whose base has maximal dimension among maps
\(S\dashrightarrow Y\) for which a multiple of \(L\) dominates an
ample base divisor in pseudoeffective order. Maximality forces ratios
of the relevant sections to be base functions. Their horizontal zero
orders consequently vary affinely with the exponent.

After birational modifications, a normalization along base divisors
produces a line bundle \(B\) on a smooth base \(Y\) and a nonzero map
\[
\sigma:f^*B\longrightarrow\ell\pi^*L,\qquad \ell>0,
\]
where \(V\) is smooth, \(\pi:V\to S\) is birational, and
\(f:V\to Y\) is a morphism. This map will be used to transfer the
sections constructed on the base to positive multiples of \(L\) on
\(S\). The same maximality gives
rank-one adjoint direct images, whose integral metrics yield a
semipositive metric on \(B\) by Berndtsson's positivity theorem
\cite{Berndtsson2009}. The normalization of \(\sigma\) makes the
resulting weights locally bounded, including near singular fibers.

These bounds allow Fujino's Koll\'ar--Nadel theorem \cite{Fujino2018}
to give vanishing at every nonnegative base twist. A single Euler
polynomial therefore equals the corresponding section dimension even
at zero, where the exceptional canonical section makes it positive.
Some positive exponent gives a section, which \(\sigma\) transfers
back to \(S\). As in the invariant-index argument, control at zero
is what rules out the zero polynomial; the two constructions of that
control are independent.

This criterion applies to differential forms through a determinant
construction. The generic-span method of Lazi\'c--Peternell
\cite[Lemma~4.1]{LP2018}, adapted in
Lazi\'c--Matsumura--Peternell--Tsakanikas--Xie
\cite[Lemma~5.1]{LMPTX2023}, extracts one fixed determinant line
from infinitely many twisted forms. We give the argument for an
unbounded positive sequence, including the numerically trivial case
excluded by the hypotheses of those two cited lemmas. The cotangent
subsheaf theorem \cite[Theorem~4.1]{LMPTX2023}, with Ou's generic
nefness result \cite[Theorem~1.4]{Ou2023} as an antecedent, makes the
negative of that line pseudoeffective. Theorem~\ref{thm:conversion}
then removes the error.

\subsection{Application through monodromy}

For a smooth connected projective \(X\), anticanonical nonvanishing asks
for a section of \(-mK_X\) for some \(m>0\). The positive-multiple
question under smooth semipositivity is associated with Yau's
anticanonical-section problem \cite[Problem~75]{Yau1994}.
The prescribed-Ricci theorem \cite{Yau1978} and the structure theorems
of Demailly--Peternell--Schneider and Campana--Demailly--Peternell
\cite{DPS1996,CDP2015} identify a compact rationally connected factor
of the universal cover. Residual monodromy can still act on it through
a compact torus, so a section on that factor alone need not descend.
M\"uller's equivariant theorem produces invariant plurisections under
semiampleness \cite[Theorem~C]{Muller2025}.
His Theorem~A proves nonvanishing for projective klt pairs with nef
anti-log-canonical divisor when its restriction to the general fiber of
the maximal rationally connected fibration is semiample; Corollary~B
proves nonvanishing for projective klt threefold pairs with nef
anti-log-canonical divisor. Our smooth application assumes metric
semipositivity and imposes no fiber-semiampleness condition.

We apply the invariant index on the compact factor \(F\) and ordinary
conversion on a finite cover \(S\) of \(X\); descending forms is the
step that connects those two spaces.

The companion \emph{Anticanonical nonvanishing from smooth
semipositivity} \cite{CompanionH} proves a finite-cover description with the
residual torus and invariant canonical frames retained, as well as the
finite \'etale norm of a section. We use only these
geometric inputs. We prove locally that invariant twisted forms on the
compact factor descend injectively to the finite cover. Theorem~\ref{thm:index}
and hard Lefschetz provide those forms, and
Theorem~\ref{thm:conversion} is applied after descent. Taking the norm
then proves that \(-mK_X\) has a nonzero section for some \(m>0\).

Passing to a positive multiple is necessary. On an Enriques surface
\(E\), \(K_E\) has order two and \(H^0(E,-K_E)=0\)
\cite{Dolgachev2016}. Consequently \(E\times\PP^1\) has a smoothly
semipositive, non-torsion anticanonical bundle without a first-power
section. In dimension two, the classification of
Chen--Filip--Sun--Tosatti--Zhang
\cite[Theorem~1.3]{CFSTZ2026} gives complementary geometric
context.\footnote{Their introduction credits ChatGPT~5.6~Sol with
suggesting a preliminary version of Lemma~5.2, which the authors
reformulated and developed into its geometric statement and proof.}

Section~\ref{sec:index} proves the index theorem and derives invariant
twisted forms. Section~\ref{sec:conversion} proves conversion by
constructing the normalized relative section and bounded base metric.
Section~\ref{sec:differentials} gives the determinant argument.
Section~\ref{sec:descent} proves form descent and the global application.
The two principal theorems are independent of one another; their
combination takes place only in that last application.
